\chapter{Kết luận và hướng phát triển}
\label{ch:concl}

\section{Kết luận}
Đề tài đã thiết kế, kiểm chứng và tích hợp một bộ tăng tốc nhân ma trận INT8 dựa trên mảng systolic weight-stationary $N\times N$ vào hệ thống RISC-V, và mở rộng thành công cho suy luận CNN (LeNet-5) và RNN (LSTM) theo yêu cầu của giảng viên. Các kết quả chính:
\begin{itemize}
  \item Mảng systolic có độ trễ $2N-1$ chu kỳ, thông lượng một hàng mỗi chu kỳ, được chứng minh bằng phân tích và quan sát trên dạng sóng; wrapper Avalon-MM bổ sung tích lũy theo $K$, bias, requant và ReLU/sigmoid/tanh.
  \item Mọi đầu ra của RTL và CPU trùng khớp từng bit với mô hình số nguyên Python: 26 unit test mảng, 24 phép thử hậu xử lý, 4 kích thước GEMM, 8 ảnh LeNet-5 và 16 ảnh LSTM cho mỗi $N\in\{4,8\}$.
  \item Độ chính xác INT8 trên \num{10000} ảnh MNIST: LeNet-5 98,85\% (float 98,83\%), LSTM 97,50\% (float 97,58\%).
  \item Với $N=8$ trên PicoRV32: LeNet-5 nhanh hơn 6,9$\times$, LSTM 3,3$\times$ (batch 1) và 16,4$\times$ (batch 16).
  \item Phân tích profile chỉ ra rằng mảng bận dưới 1\% thời gian; nút thắt là im2col/pooling trên CPU và chép dữ liệu qua MMIO.
  \item Trên Cyclone~V (Quartus Lite 25.1std), \code{npu\_avmm} đứng riêng fit với mọi cấu hình. RTL gốc ($N=8$) dùng \num{18521} ALM (44\%) và chỉ đạt 52\,MHz vì bộ đệm bị tổng hợp thành flip-flop. Biến thể \code{rtl\_fpga} đặt bộ đệm vào M10K và đường ống hóa requant, cho kết quả bit-exact như bản gốc, chỉ dùng \num{2187} ALM (5\%) và 48 DSP, đạt 127,5\,MHz. Luồng mở Yosys/nextpnr trên ECP5 xác nhận cùng nguyên nhân (bộ đệm chiếm khoảng 90\% diện tích), nhưng không cấu hình nào của RTL gốc đạt 50\,MHz.
  \item Luồng hiện thực trên DE10-Nano với Nios~V đã được dựng và lớp nền tảng đã được kiểm tra trên đúng bản đồ bộ nhớ của board; \todo{bổ sung tài nguyên toàn hệ thống và số đo trên board khi có}.
\end{itemize}

\section{Hướng phát triển}
Các hướng dưới đây được sắp theo lợi ích kỳ vọng, dựa trực tiếp vào phân tích nút thắt ở Chương~\ref{ch:results}:
\begin{enumerate}
  \item \textbf{DMA (Avalon-MM master)}: cho \NPU{} tự đọc trọng số/kích hoạt và ghi kết quả từ RAM theo bộ mô tả (descriptor), giải phóng CPU khỏi khoảng 30\% thời gian chép dữ liệu.
  \item \textbf{im2col trong phần cứng}: một bộ sinh địa chỉ đọc cửa sổ $k\times k$ trực tiếp từ bản đồ đặc trưng, kết hợp với DMA; cùng với \textbf{maxpool phần cứng} ở đầu ra QBUF, điều này loại bỏ phần ``khác'' 63\% và phá vỡ giới hạn Amdahl 10,9$\times$.
  \item \textbf{Double buffering} ABUF/trọng số và \textbf{ngắt} thay polling: nạp khối kế tiếp trong khi mảng tính khối hiện tại.
  \item \textbf{Bộ đệm lớn hơn} (M10K): giữ nhiều khối trọng số và toàn bộ bản đồ đặc trưng trên chip; tận dụng tái sử dụng trọng số khi $M$ lớn.
  \item \textbf{Hỗ trợ LSTM tốt hơn}: thực hiện cập nhật element-wise $c_t$, $h_t$ trong phần cứng; thử $H$ lớn hơn (64, 128) và xếp lịch nhiều chuỗi để lấp đầy mảng khi không thể batching.
  \item \textbf{Khai thác độ thưa (sparsity)}: bỏ qua các khối trọng số bằng 0 hoặc kích hoạt bằng 0 sau ReLU.
  \item \textbf{Lượng tử hóa theo kênh} (một cặp $(\mu,\sigma)$ cho mỗi cột) và zero-point bất đối xứng, cho các mạng lớn hơn MNIST.
  \item \textbf{Tích hợp chặt hơn với RISC-V}: lệnh tùy biến (kiểu RoCC của Gemmini \cite{genc2021gemmini}) hoặc giao diện coprocessor để giảm chi phí điều khiển mỗi lượt.
  \item \textbf{Kiểm chứng mạnh hơn}: assertion SVA, kiểm chứng hình thức cho \code{sa\_top} (độ trễ, thứ tự hàng), đo coverage.
  \item \textbf{Luồng ASIC}: tổng hợp và đặt-định tuyến với một PDK mở (ví dụ SKY130 với OpenROAD/OpenLane) để đo diện tích, công suất và tần số; so sánh năng lượng mỗi phép MAC với CPU.
\end{enumerate}
